%==============================================================================
% jarvis-preamble.tex
% Shared packages, metadata macros, and semantic text commands for the
% Jarvis-HEP manual.  Loaded by main.tex after \documentclass{tufte-book}.
%
% The tufte-book class already provides: xcolor, hyperref (unless nohyper),
% geometry, graphicx, \allcaps, \smallcaps, \newthought, \sidenote,
% \marginnote, \footnote->sidenote, fullwidth env, marginfigure, etc.
% We only add what a code/reference manual additionally needs.
%==============================================================================

% --- Math ---------------------------------------------------------------------
\usepackage{amsmath}
\usepackage{amssymb}

% --- Tables -------------------------------------------------------------------
\usepackage{booktabs}
\usepackage{longtable}
\usepackage{array}
\usepackage{ragged2e}      % \RaggedRight inside narrow table cells

% --- Misc ---------------------------------------------------------------------
\usepackage{pmboxdraw}     % render block/box-drawing glyphs (the Jarvis -v banner) in pdfLaTeX
\usepackage{newunicodechar} % map the literal dot glyph U+2B24 (the banner icon) to a drawn circle
\usepackage{varwidth}      % shrink-to-content box so the banner panel hugs its text (no empty margin)
\usepackage{xspace}        % smart trailing space for text macros
\usepackage{enumitem}      % tighter description/itemize lists
\setlist{noitemsep, topsep=2pt}

% --- Project version single-source --------------------------------------------
% This manual is a LIVING document: it is revised in step with each Jarvis-HEP
% release.  \JversionNum names the release this particular build documents --
% bump it (and rebuild) on every release; do not hard-code version numbers in
% the prose, always use \Jversion so there is a single source of truth.
\newcommand{\JversionNum}{1.7.3}
\newcommand{\Jversion}{v\JversionNum\xspace}

% --- Product / ecosystem names ------------------------------------------------
\newcommand{\Jarvis}{\textsc{Jarvis-HEP}\xspace}
\newcommand{\Operas}{\textsc{Jarvis-Operas}\xspace}
\newcommand{\JPlot}{\textsc{Jarvis-PLOT}\xspace}

% --- Inline semantic markup ---------------------------------------------------
% Monospace helper that lets long tokens break in the narrow text column: an
% underscore (\_) and a dot offer a zero-penalty breakpoint, so names like
% interpolations\_1D or Runtime.batch\_size can wrap instead of overflowing.
\newcommand{\jbttbreak}{\penalty0\hskip0pt\relax}
\newcommand{\jtt}[1]{\begingroup
  \def\_{\textunderscore\jbttbreak}%
  \texttt{#1}\endgroup}
% Code-ish things rendered in monospace with consistent meaning across the book.
\newcommand{\code}[1]{\jtt{#1}}                          % generic inline code
\newcommand{\yamlkey}[1]{\jtt{#1}}                       % a YAML key / block name
\newcommand{\cli}[1]{\jtt{#1}}                           % a shell / CLI fragment
\newcommand{\token}[1]{\jtt{#1}}                         % runtime token @SampleID etc.
\newcommand{\marker}[1]{\jtt{#1}}                        % path marker &J/
\newcommand{\file}[1]{\jtt{#1}}                          % a file or path (overridden in dirtree)
\newcommand{\sampler}[1]{\textsf{#1}}                    % a sampler method name
\newcommand{\optitem}[1]{\textit{#1}}                    % an "(optional)" tag etc.

% Reference to a sampler JSON schema file (authoritative key source).
\newcommand{\schemaref}[1]{\texttt{card/schema/#1}}

% A "registered alias" inline note, e.g. \alias{ToyMCMC}{MCMC}
\newcommand{\alias}[2]{\texttt{#1}~(alias of~\texttt{#2})}

% --- Required / Optional inline tags for key tables ---------------------------
\newcommand{\req}{\textbf{required}}
\newcommand{\opt}{optional}

% --- Tufte: \paragraph is our de-facto third heading level --------------------
% tufte-book forbids \subsubsection, so \paragraph (provided by the class) is the
% deepest heading we use.  We deliberately do NOT redefine it, to avoid clashing
% with the class's titlesec-based sectioning.

% --- A compact key/value description list for YAML key references -------------
\newenvironment{keylist}
  {\begin{description}[style=nextline, leftmargin=1.2em, font=\ttfamily]}
  {\end{description}}

% --- Coloured cross-reference & citation links --------------------------------
% tufte loads hyperref with colorlinks=false (links are black, box-less), so
% references look like plain text.  Turn on coloured links so references to
% figures, tables, equations, and sections -- and citations -- stand out.
% (hyperref is already loaded by the tufte-book class at this point.)
\definecolor{jlink}{HTML}{1F6FB2}  % cross-refs: figures, tables, equations, sections
\definecolor{jcite}{HTML}{2E7D32}  % citations / bibliography
\definecolor{jurl}{HTML}{1A5FB4}   % URLs
\hypersetup{
  colorlinks = true,
  linkcolor  = jlink,
  citecolor  = jcite,
  urlcolor   = jurl,
  linktoc    = page,   % keep ToC titles black; colour only the page-number links
}
